\section*{Abstract}
% \addcontentsline{toc}{section}{Abstract}
% \blindtext

There often appear cases where a particular marginal distribution in the generative process of a diffusion model is undesirable or better replaced with another distribution to suit a specific task. Scenarios such as fair generation or label shift are common examples of this need. Using Jeffrey's Rule, it is possible to transform this undesired distribution into a more suitable target distribution. In the context of distribution-guided diffusion, actually sampling from the posterior requires an accurate estimate of the clean data $x_0$ from a noisy sample $x_t$. Tweedie's formula is often employed for this purpose but its use of a single expected value, although good for simple classifier-guided diffusion, may not be best suited for a more complex case such as distribution-guidance. This thesis explores a Markov Chain Monte Carlo (MCMC) approach for sampling from the marginal distribution defined by Jeffreys Rule, investigating its performance within this distribution-guided diffusion setting.
VERSION 2:

While diffusion models capture the data distributions they are trained on, the resulting learned marginals may require modification to satisfy constraints such as fairness or to account for label shift. Jeffrey's rule provides a principled mechanism for transforming an undesired marginal distribution into a desired target distribution.

Sampling from the resulting posterior during the diffusion process requires estimating the clean data $x_0$ from a noisy observation $x_t$. A common approach is to use Tweedie's formula, which provides an estimate of $x_0$ through the posterior mean. While effective in many settings, this mean-based approximation may be insufficient when the target distribution is complex or multi-modal, as can occur in distribution-guided diffusion models.

This thesis investigates the use of Markov Chain Monte Carlo (MCMC) methods to sample from the marginal distribution induced by Jeffrey's rule, evaluating their effectiveness within a distribution-guided diffusion framework.
